\subsection{End-to-End Evaluation and Comparison}
\label{sec:endtoend}

We evaluate the full pipeline with the default components, \emph{Qwen3.8-2.4T-A95B} with \emph{Examples} as classifier and \emph{Claude Sonnet 5.5} with \emph{Multi-Step} as mapper, in two variants: classifier plus mapper (C+MS) and the mapper alone, which relies on its rule filter (Multi-Step). Each variant runs three times, and Table~\ref{tab:main_results} reports the mean. Baseline results on Hospital, Restaurant, SISU, and Tax come from the preliminary study (Section~\ref{sec:implementation}), with the same data and gold constraints (their SISU violation scores predate the pair semantics of Section~\ref{sec:mapper:sql}); the baselines were not run on the HMDA datasets.

\begin{table}[!ht]
  \centering
  \caption{End-to-end comparison of \system\ (C+MS and Multi-Step; mean of three runs) and the baselines. \textbf{Profiling}: DC discovery from data alone; \textbf{Text2SQL}: XiYan-SQL adapted to generate violation queries; -- marks configurations not run; n/a, metric not applicable; bold, best F$_1$ per metric and dataset.}
  \label{tab:main_results}
  \small
  \setlength{\tabcolsep}{4pt}
  \resizebox{\columnwidth}{!}{%
  \begin{tabular}{llccccccccc}
\toprule
\textbf{Family} & \textbf{Method} & \multicolumn{3}{c}{Constr. Classification} & \multicolumn{3}{c}{Constr. Extraction} & \multicolumn{3}{c}{Viol. Detection} \\
\cmidrule(lr){3-5}  \cmidrule(lr){6-8}  \cmidrule(lr){9-11}
 &  & P & R & F$_1$ & P & R & F$_1$ & P & R & F$_1$ \\
\midrule
\addlinespace[2pt]
\multicolumn{11}{c}{\textbf{Hospital}} \\
\addlinespace[1pt]
\multirow{2}{*}{\system} & ~~C+MS & 0.81 & 1.00 & \textbf{0.89} & 1.00 & 1.00 & \textbf{1.00} & 1.00 & 1.00 & \textbf{1.00} \\
 & ~~Multi-Step & 0.48 & 1.00 & 0.65 & 1.00 & 1.00 & \textbf{1.00} & 1.00 & 1.00 & \textbf{1.00} \\
\cmidrule(l){2-11}
\multirow{1}{*}{Profiling} & ~~DC Discovery & n/a & n/a & n/a & 0.00 & 0.00 & 0.00 & 0.83 & 1.00 & 0.91 \\
\cmidrule(l){2-11}
\multirow{1}{*}{Text2SQL} & ~~XiYan-SQL & 0.50 & 1.00 & 0.67 & 0.60 & 1.00 & 0.75 & 0.83 & 1.00 & 0.91 \\
\midrule
\addlinespace[2pt]
\multicolumn{11}{c}{\textbf{Restaurant}} \\
\addlinespace[1pt]
\multirow{2}{*}{\system} & ~~C+MS & 0.89 & 0.71 & 0.79 & 0.80 & 0.57 & \textbf{0.67} & 1.00 & 0.34 & 0.35 \\
 & ~~Multi-Step & 0.65 & 0.71 & 0.68 & 0.80 & 0.57 & \textbf{0.67} & 0.87 & 0.01 & 0.02 \\
\cmidrule(l){2-11}
\multirow{1}{*}{Profiling} & ~~DC Discovery & n/a & n/a & n/a & 0.00 & 0.00 & 0.00 & 0.77 & 0.87 & \textbf{0.82} \\
\cmidrule(l){2-11}
\multirow{1}{*}{Text2SQL} & ~~XiYan-SQL & 0.86 & 0.75 & \textbf{0.80} & 0.17 & 0.12 & 0.14 & 0.45 & 0.64 & 0.52 \\
\midrule
\addlinespace[2pt]
\multicolumn{11}{c}{\textbf{SISU}} \\
\addlinespace[1pt]
\multirow{2}{*}{\system} & ~~C+MS & 0.67 & 1.00 & \textbf{0.80} & 0.65 & 0.65 & \textbf{0.65} & 0.03 & 0.98 & \textbf{0.06} \\
 & ~~Multi-Step & 0.65 & 1.00 & 0.79 & 0.65 & 0.65 & \textbf{0.65} & 0.03 & 0.97 & 0.05 \\
\cmidrule(l){2-11}
\multirow{1}{*}{Profiling} & ~~DC Discovery & n/a & n/a & n/a & 0.00 & 0.00 & 0.00 & 0.00 & 0.99 & 0.01 \\
\cmidrule(l){2-11}
\multirow{1}{*}{Text2SQL} & ~~XiYan-SQL & 1.00 & 0.65 & 0.79 & 0.45 & 0.29 & 0.36 & 0.00 & 1.00 & 0.01 \\
\midrule
\addlinespace[2pt]
\multicolumn{11}{c}{\textbf{Tax}} \\
\addlinespace[1pt]
\multirow{2}{*}{\system} & ~~C+MS & 0.71 & 1.00 & \textbf{0.82} & 1.00 & 1.00 & \textbf{1.00} & 0.59 & 1.00 & \textbf{0.70} \\
 & ~~Multi-Step & 0.33 & 1.00 & 0.50 & 1.00 & 1.00 & \textbf{1.00} & 0.01 & 1.00 & 0.02 \\
\cmidrule(l){2-11}
\multirow{1}{*}{Profiling} & ~~DC Discovery & n/a & n/a & n/a & 0.00 & 0.00 & 0.00 & 0.00 & 0.94 & 0.01 \\
\cmidrule(l){2-11}
\multirow{1}{*}{Text2SQL} & ~~XiYan-SQL & 0.40 & 1.00 & 0.57 & 0.29 & 0.67 & 0.40 & 0.01 & 1.00 & 0.02 \\
\midrule
\addlinespace[2pt]
\multicolumn{11}{c}{\textbf{HMDAZ}} \\
\addlinespace[1pt]
\multirow{2}{*}{\system} & ~~C+MS & 0.86 & 0.89 & 0.87 & 0.13 & 0.15 & 0.14 & 0.26 & 1.00 & 0.38 \\
 & ~~Multi-Step & 0.96 & 0.89 & \textbf{0.92} & 0.25 & 0.26 & \textbf{0.25} & 0.48 & 1.00 & \textbf{0.64} \\
\cmidrule(l){2-11}
\multirow{1}{*}{Profiling} & ~~DC Discovery & -- & -- & -- & -- & -- & -- & -- & -- & -- \\
\cmidrule(l){2-11}
\multirow{1}{*}{Text2SQL} & ~~XiYan-SQL & -- & -- & -- & -- & -- & -- & -- & -- & -- \\
\midrule
\addlinespace[2pt]
\multicolumn{11}{c}{\textbf{HMDA Checks}} \\
\addlinespace[1pt]
\multirow{2}{*}{\system} & ~~C+MS & 0.90 & 0.93 & 0.91 & 0.45 & 0.43 & 0.44 & 0.99 & 0.56 & \textbf{0.71} \\
 & ~~Multi-Step & 0.90 & 0.93 & \textbf{0.92} & 0.45 & 0.44 & \textbf{0.45} & 0.99 & 0.56 & \textbf{0.71} \\
\cmidrule(l){2-11}
\multirow{1}{*}{Profiling} & ~~DC Discovery & -- & -- & -- & -- & -- & -- & -- & -- & -- \\
\cmidrule(l){2-11}
\multirow{1}{*}{Text2SQL} & ~~XiYan-SQL & -- & -- & -- & -- & -- & -- & -- & -- & -- \\
\bottomrule
\end{tabular}

  }%

\end{table}

\system\ is the only approach that produces gold constraints on every dataset. Extraction F$_1$ is 1.00 on Hospital and Tax, 0.67 on Restaurant, 0.65 on SISU, 0.44--0.45 on HMDA Checks, and 0.14--0.25 on HMDAZ (0.35 with equivalence-aware matching), against at most 0.75 for XiYan-SQL and zero for DC discovery. Classification F$_1$ of C+MS lies between 0.79 and 0.91, including on the wide HMDA schemas.

The classifier mainly improves precision: compared with Multi-Step alone, C+MS raises classification precision on Hospital (0.81 vs.\ 0.48), Restaurant (0.89 vs.\ 0.65), and Tax (0.71 vs.\ 0.33). On Tax, this decides the aggregate violation score. Multi-Step alone maps about 12 non-constraint sentences per run, written to resemble constraints; one turns ``major cities exist within fixed states'' into \texttt{t1.city = t2.city AND t1.state != t2.state}, which flags 52k tuples, so the aggregate falls to 0.02. C+MS filters most of them and reaches 0.70 on average (0.55--1.00). On the HMDA datasets, both variants make almost the same decisions; on HMDAZ, they still write different DCs for the same triggers (exact extraction 0.14 vs.\ 0.25, equivalence-aware 0.35 for both), and the aggregate varies from 0.02 to 0.75 across runs.

The violation results expose two further failure modes. First, a broad DC on a non-constraint sentence can dominate the aggregate: on SISU, both variants map sentences that define quota categories, labeled as non-constraints, to checks that flag up to 12k tuples, so the aggregate stays at 0.05--0.06 although the gold rules are recovered well (per-rule violation F$_1$ 0.91--0.95, execution equivalence 0.86--0.90). Second, missing a dominant rule: on Restaurant, the yearly-inspection rule flags 99\% of the gold violations, and Multi-Step alone returns no DC for it in all three runs, C+MS in two of them, because the rule filter judges this cross-year requirement as not expressible, although the derived attribute encodes it. The aggregate is therefore 0.02 for Multi-Step and 0.35 for C+MS, below DC discovery (0.82). Both failures depend on the backbone and the run: with GPT-6 Sol as the mapper, the pipeline reached 1.00 on Restaurant and up to 0.98 on SISU, but 0.15 on Tax. On HMDA Checks, the aggregate of 0.71 is not evidence of quality: no date-range edit is mapped, and the score comes from an incorrect DC that flags 481k tuples, most of which violate some gold rule. There, 83\% of the gold rules with violations are reproduced exactly, and the mean per-rule F$_1$ is 0.77--0.78.

The baselines are competitive only on narrower parts of the experiments. DC discovery recovers violation signals when the data distribution strongly reflects the target rules, as on Hospital and Restaurant, but never the gold constraints (extraction F$_1$ zero on all datasets). XiYan-SQL is competitive on Hospital, but on SISU and Tax it returns overly broad queries, with high violation recall and near-zero precision. Since a single gold rule flags most gold-violating tuples in four of the six datasets, broad violation sets can still overlap strongly with the gold ones, while a correct pipeline can be penalized by one broad false positive; we therefore recommend reading the aggregate together with the per-rule metrics. Overall, \system\ is the only approach that remains strong on classification and extraction across all six datasets; it is behind the baselines only on Restaurant, in violation detection and, by 0.01, in classification.

\smallskip\noindent\textbf{Usability and cost.}
The returned set stays small and reviewable: C+MS returns 6--25 DCs per run on Hospital, Restaurant, SISU, and Tax, 9 on HMDAZ, and 95 on HMDA Checks, close to the number of gold constraints (6--92). Users can rank the DCs by the number of tuples they flag, read the source sentence of each, and use the optional review step (Section~\ref{sec:overview}); the broad DCs above stand out precisely by their violation counts.
All reported experiments cost US\$166 through OpenRouter. With Claude Sonnet 5.5, an end-to-end run costs US\$0.4--2.9 on Hospital, Restaurant, SISU, and Tax and about US\$9.7 on HMDA Checks (158 sentences, 1M rows). The classifier makes the pipeline up to six times cheaper (Tax), since rejected sentences never reach the mapper; on HMDA Checks, where most sentences are constraints, both variants cost the same.
